% Notes of the Montreal seminar, 1970 (cote 129, typed part): archive cote 129, site page 75 (PDF page 76).
% typescript. Draft compared: none.
% AI-transcribed by Claude Opus 5.5 (claude-opus-5-5, adaptive thinking, effort high), Message Batches API, batch mtl129_full (msgbatch_01AGnDKHwy6ATCM7ehUa6263), 2026-11-13; not yet checked by eye.
\scanpage{129/75}
\mspage{\hand{III-20}}

p, où l'on a des points isolés $s$ où l'invariant de Hasse de $A_s$ est
nul. On a cependant le résultat suivant:

7.4. \uline{Proposition}. Soit $G$ un groupe de Barsotti-Tate sur $S$. Les
propriétés suivantes sont équivalentes:

i) $G$ est extension d'un groupe de BT ind-étale par un groupe
de BT à fibres connexes.

ii) Pour tout $n \in \mathbb{N}$, $G(n)$ est extension d'un groupe fini
étale par un groupe fini infinitésimal.

ii)$^{\text{bis}}$ $G(1)$ est extension d'un groupe fini étale par un groupe
fini infinitésimal.

iii) Pour tout $n \in \mathbb{N}$, la fonction $s \mapsto \text{rang sep } G(n)_s$ est lo-%
calement constante.

iii)$^{\text{bis}}$ La fonction $n \mapsto \text{rang sep } G(1)_s$ est localement constante.

L'équivalence de (iii) et (iii)$^{\text{bis}}$ résulte de ce que le rang sé-%
parable de $G(n)_s$ est la puissance n-ième du rang séparable de $G(\unsure{i|1})_s$ .
D'autre part, (i) résulte immédiatement de (ii), et, d'après ce qui précède,
la seule implication à démontrer est (iii) $\Rightarrow$ (ii). On utilise pour cela
le lemme suivant (EGA, IV, ):

\uline{Lemme}. Soit $f : X \to S$ un morphisme fini localement libre de schémas.
Pour que $f$ se factorise en $f'$ étale et $f''$ radiciel et surjectif, il
\diagram{\begin{tikzcd}[column sep=small]
X \arrow[rr, "f''"] \arrow[dr, "f"'] & & Y \arrow[dl, "f'"] \\
& S &
\end{tikzcd}}
faut et il suffit que la fonction $s \mapsto \text{rang sep } (X_s)$
soit localement constante. La factorisation est alors
unique et fonctorielle en $f$.

(De la fonctorialité, il résulte que si $X$ est un schéma en groupes sur
$S$, $Y$ l'est aussi.)

\margin{75}
